\documentclass[11pt,a4paper]{article}

% ---------------------------------------------------------
% PREAMBLE (stable, warning-free, Overleaf-safe)
% ---------------------------------------------------------
\usepackage[utf8]{inputenc}
\usepackage[T1]{fontenc}
\usepackage{lmodern}
\usepackage{amsmath, amssymb, amsfonts}
\usepackage{bm}
\usepackage{geometry}
\usepackage{graphicx}
\usepackage{hyperref}
\usepackage{cite}
\usepackage{authblk}
\usepackage{physics}
\usepackage{mathtools}

\geometry{margin=2.4cm}

\title{\textbf{Companion LIII: E--B Mode Conversion in the Relaxation-Driven Cyclic Cosmology}}
\author{Michael Lehmann \\ \texttt{mi.lehmann@gmx.de}}
\date{April 2026}

\begin{document}
\maketitle

\begin{abstract}
Weak gravitational lensing converts CMB E-modes into B-modes, providing a powerful 
probe of the gravitational potential along the line of sight. In the standard 
cosmological model, this conversion is determined by the lensing potential and the 
geometry of the deflection field. In the Relaxation-Driven Cyclic Cosmology (RDCC), 
the scale-dependent evolution of the gravitational potential modifies the lensing 
kernel, the deflection field, and the E--B mode coupling in a characteristic way. 
This Companion develops a continuous description of E--B mode conversion in RDCC, 
deriving how the relaxation scale influences the lensing-induced B-mode spectrum, 
the mode-coupling structure, and the coherence of the deflection field. The 
resulting framework provides a distinctive observational signature of the RDCC 
gravitational field.
\end{abstract}
% ---------------------------------------------------------
\section{Introduction}

CMB polarization provides a unique window into the geometry and dynamics of the 
Universe. E-modes arise from scalar perturbations, while B-modes are generated by 
tensor perturbations or by the lensing-induced conversion of E-modes into B-modes. 
The latter process is sensitive to the gravitational potential along the line of 
sight and provides a direct probe of large-scale structure.

In the standard cosmological model, the E--B mode conversion is determined by the 
lensing potential and the geometry of the deflection field. In RDCC, the 
gravitational potential evolves differently. The relaxation mechanism suppresses 
the decay of small-scale modes and enhances the coherence of large-scale modes. 
This modifies the lensing potential, the deflection field, and the E--B mode 
coupling in a scale-dependent manner, producing a distinctive signature in the 
lensing-induced B-mode spectrum.
% ---------------------------------------------------------
\section{The RDCC Lensing Potential and Deflection Field}

The lensing potential is given by
\begin{equation}
    \phi(\hat{\mathbf{n}})
    = -2 \int_{0}^{\chi_{*}}
      d\chi\,
      \frac{\chi_{*} - \chi}{\chi_{*}\chi}\,
      \left[\Phi + \Psi\right](\chi,\hat{\mathbf{n}}).
\end{equation}

In RDCC, the metric potentials evolve as
\begin{align}
    \Phi_{\mathrm{RDCC}}(k,a)
    &= \Phi(k,a)
       \left[
         1 - \chi_{\Phi}
         \left(\frac{k}{k_{\mathrm{rel}}}\right)^{\alpha}
       \right], \\
    \Psi_{\mathrm{RDCC}}(k,a)
    &= \Psi(k,a)
       \left[
         1 - \chi_{\Psi}
         \left(\frac{k}{k_{\mathrm{rel}}}\right)^{\alpha}
       \right].
\end{align}

The deflection field becomes
\begin{equation}
    \mathbf{d}_{\mathrm{RDCC}}
    = \nabla \phi_{\mathrm{RDCC}}
    = \mathbf{d}
      \left[
        1 - \chi_{d}
        \left(\frac{k}{k_{\mathrm{rel}}}\right)^{\alpha}
      \right].
\end{equation}
% ---------------------------------------------------------
\section{E--B Mode Conversion in RDCC}

Lensing converts E-modes into B-modes according to
\begin{equation}
    B_{\ell m}
    = \sum_{\ell' m'}
      \sum_{LM}
      \phi_{LM}\,
      \mathcal{G}_{\ell \ell' L}^{m m' M}\,
      E_{\ell' m'},
\end{equation}
where $\mathcal{G}$ is a geometric coupling coefficient.

In RDCC, the modified lensing potential produces
\begin{equation}
    B_{\ell m}^{\mathrm{RDCC}}
    = B_{\ell m}
      \left[
        1 - \chi_{B}
        \left(\frac{\ell}{\ell_{\mathrm{rel}}}\right)^{\alpha}
      \right].
\end{equation}

This modification suppresses small-scale B-modes and enhances large-scale 
coherence.
% ---------------------------------------------------------
\section{The RDCC B-Mode Power Spectrum}

The lensing-induced B-mode power spectrum is given by
\begin{equation}
    C_{\ell}^{BB}
    = \int d^{2}\ell'\,
      \left[
        \ell' \cdot (\ell - \ell')
      \right]^{2}
      C_{\ell'}^{EE}
      C_{|\ell - \ell'|}^{\phi\phi}.
\end{equation}

In RDCC, the lensing potential power spectrum becomes
\begin{equation}
    C_{L}^{\phi\phi,\mathrm{RDCC}}
    = C_{L}^{\phi\phi}
      \left[
        1 - \chi_{\phi\phi}
        \left(\frac{L}{L_{\mathrm{rel}}}\right)^{\alpha}
      \right].
\end{equation}

The B-mode spectrum becomes
\begin{equation}
    C_{\ell}^{BB,\mathrm{RDCC}}
    = C_{\ell}^{BB}
      \left[
        1 - \chi_{BB}
        \left(\frac{\ell}{\ell_{\mathrm{rel}}}\right)^{\alpha}
      \right].
\end{equation}
% ---------------------------------------------------------
\section{Large-Scale Coherence and RDCC Signatures}

The relaxation mechanism enhances the coherence of large-scale modes. The 
lensing-induced B-mode spectrum therefore exhibits:

\begin{itemize}
    \item enhanced power at large angular scales,
    \item suppressed power at small scales,
    \item a characteristic turnover near $\ell_{\mathrm{rel}}$,
    \item a modified phase structure in E--B coupling.
\end{itemize}

These signatures provide a direct observational probe of the RDCC gravitational 
field.
% ---------------------------------------------------------
\section{Conclusion}

The E--B mode conversion in the Relaxation-Driven Cyclic Cosmology reflects the 
scale-dependent evolution of the gravitational potential. The relaxation mechanism 
modifies the lensing potential, the deflection field, and the E--B mode coupling in 
a scale-dependent manner, producing distinctive signatures in the lensing-induced 
B-mode spectrum. These effects provide a direct observational window into the 
relaxation scale and the temporal structure of the RDCC gravitational field. The 
present Companion establishes the theoretical foundation for future studies of CMB 
polarization, lensing reconstruction, and the interplay between matter and gravity 
in RDCC.

\bibliographystyle{apsrev4-2}
\nocite{*}
\bibliography{references}

\end{document}
